\documentclass[11pt]{article}

\usepackage[a4paper, margin=1in]{geometry}
\usepackage[utf8]{inputenc}
\usepackage{graphicx}
\usepackage{natbib}
\usepackage{booktabs}
\usepackage{hyperref}

\title{Framing effects on consumer purchase intention: \\
       a mediation and moderation analysis}
\author{Your Name}
\date{\today}

\begin{document}
\maketitle

\begin{abstract}
We test whether messaging frames (scarcity, sustainability, social proof,
and a combined frame) increase purchase intention for everyday consumer
products, whether this effect is mediated by perceived value, and whether
it is moderated by price sensitivity. Data, scripts, and outputs live in
the same git repository as this manuscript.
\end{abstract}

\section{Introduction}

Mediation and moderation logic follows the classical framework
\citep{baron1986} and the regression-based approach in \citet{hayes2017}.

\section{Methods}

Three hundred simulated participants were randomly assigned to one of five
conditions and reported their attitudes and a hypothetical purchase
decision. See the project repository for the full codebook.

\section{Results}

\subsection{Main effect}

\begin{figure}[h]
  \centering
  \includegraphics[width=0.85\textwidth]{../figures/fig_01_pi_by_condition.png}
  \caption{Purchase intention by condition. Red points are condition means
           with 95\% confidence intervals.}
  \label{fig:main}
\end{figure}

\subsection{Mediation through perceived value}

\begin{figure}[h]
  \centering
  \includegraphics[width=0.85\textwidth]{../figures/fig_02_mediation_scatter.png}
  \caption{Perceived value predicts purchase intention; the conditions
           shift perceived value upward, producing an indirect path.}
  \label{fig:mediation}
\end{figure}

\subsection{Moderation by price sensitivity}

\begin{figure}[h]
  \centering
  \includegraphics[width=0.85\textwidth]{../figures/fig_03_moderation.png}
  \caption{The condition boost is attenuated for high-price-sensitive
           respondents.}
  \label{fig:moderation}
\end{figure}

\section{Discussion}

The figures referenced above are produced by \texttt{R/03\_analysis.R}.
Re-running that script overwrites them; this manuscript picks up the new
versions on the next compile. No file is duplicated between code and paper.

\bibliographystyle{plain}
\bibliography{../references/references}

\end{document}
